\documentclass[12pt,a4paper]{article}

% --- Packages ---
\usepackage[T1]{fontenc}
\usepackage[utf8]{inputenc}
\usepackage[english]{babel}
\usepackage{amsmath, amssymb}
\usepackage{graphicx}
\usepackage{booktabs}
\usepackage{caption}
\usepackage{subcaption}
\usepackage{setspace}
\usepackage{geometry}
\usepackage{hyperref}
\usepackage{natbib}
\usepackage{float}
\usepackage{lmodern}

\geometry{margin=2.5cm}
\onehalfspacing
\hypersetup{colorlinks=true, linkcolor=blue, citecolor=blue, urlcolor=blue}

% --- Title ---
\title{\textbf{Does Public Education Expenditure Reduce Income Inequality?\\
Cross-Country Evidence from the World Bank}}

\author{
  Jane Doe\thanks{École Polytechnique -- Institut Polytechnique de Paris.
  Email: \texttt{jane.doe@polytechnique.edu}} \and
  John Smith\thanks{École Polytechnique -- Institut Polytechnique de Paris.
  Email: \texttt{john.smith@polytechnique.edu}}
}

\date{May 2026}

\begin{document}

\maketitle

\begin{abstract}
  \noindent
  We study the relationship between government expenditure on education and income
  inequality across a panel of 10 countries over the period 2015--2019, using data from
  the World Bank World Development Indicators. Exploiting within-country variation in
  education spending and controlling for GDP per capita, secondary school enrollment,
  and urbanization, we find a negative and statistically significant association between
  education expenditure and the Gini index: a one percentage-point increase in education
  spending as a share of GDP is associated with a reduction of approximately 1.5 Gini
  points. Country fixed-effects estimates confirm this pattern, though the effect is
  attenuated when year fixed effects are included. These results are consistent with
  the hypothesis that public education investment mitigates inequality by expanding
  access to human capital accumulation.\\[4pt]
  \textbf{JEL Codes:} I24, D63, H52\\
  \textbf{Keywords:} income inequality, education expenditure, human capital, Gini index
\end{abstract}

\newpage

% =============================================================================
\section{Introduction}
\label{sec:intro}
% =============================================================================

\subsection*{Motivation}

Income inequality has increased substantially in many advanced and developing economies
over the past four decades \citep{piketty2014capital, atkinson2015inequality}. High
levels of inequality carry significant consequences for economic growth, social mobility,
and political stability. Understanding the determinants of inequality---and identifying
effective policy levers to address it---has therefore become a central question in
economics and public policy. Among the many instruments available to governments,
education policy occupies a special place: it can both directly redistribute income
(by improving the earning capacity of lower-income households) and indirectly reduce
inequality through its effects on social mobility and the skill composition of the labour
force \citep{goldin2008race}.

Despite its centrality in policy discussions, the empirical evidence on the effect of
\textit{public} education expenditure on income inequality remains mixed. Cross-country
comparisons are complicated by differences in the composition of spending (primary vs.
tertiary), institutional quality, and the level of development. In addition, reverse
causality is a genuine concern: governments facing high inequality may have greater
incentives to invest in education, creating upward bias in naïve OLS estimates.

\subsection*{Research Question}

This paper asks: \textit{Does higher government expenditure on education reduce income
inequality?} We study this question using a cross-country panel from the World Bank
World Development Indicators, focusing on 10 heterogeneous countries over 2015--2019.

\subsection*{Brief Summary}

We use annual data on the Gini index, government education expenditure as a share of
GDP, GDP per capita, secondary school enrollment, and urbanization for 10 countries
spanning different income levels and regions. Our identification strategy exploits
within-country variation in education spending over time, relying on country fixed
effects to absorb time-invariant confounders such as historical institutions and
culture. Baseline OLS estimates suggest a strong negative correlation between education
spending and the Gini index; this finding is robust to the inclusion of country fixed
effects, although attenuated in magnitude. We interpret these results as consistent
with the human capital channel: education expenditure expands access to schooling among
lower-income households, compressing the skill and wage distribution.

\subsection*{Literature Review}

Our paper relates to two strands of the literature. First, a body of work studies the
role of education in shaping long-run inequality. \citet{goldin2008race} argue that
the US inequality surge of the 1980s reflects a ``race between education and technology''
in which technological change outpaced the supply of skilled workers; their framework
implies that policies expanding educational attainment should reduce inequality when
labour supply responds sufficiently. \citet{checchi2003education} provides cross-country
evidence that education affects income distribution through both wage compression and
intergenerational mobility, though the direction depends on how education is financed
and distributed across the income spectrum.

Second, we contribute to the fiscal policy and inequality literature. \citet{clements2012}
and \citet{imf2014fiscal} document that social spending---including education---tends to
reduce market income inequality in OECD countries, while the effect is weaker in
developing economies, partly because spending often disproportionately benefits the
middle class rather than the poor. Our results corroborate these findings but highlight
the importance of controlling for GDP per capita and urbanization when comparing
countries at different development stages.

% =============================================================================
\section{Descriptive Evidence}
\label{sec:data}
% =============================================================================

\subsection{Data Description}

Our dataset is constructed from the \textbf{World Bank World Development Indicators}
(WDI), one of the most widely used sources of international development data. The WDI
compiles data from officially recognised international sources and covers over 200
countries and territories. We use five indicators for the period 2015--2019:
\begin{itemize}
  \item \textbf{Gini index} (\texttt{SI.POV.GINI}): a measure of income distribution
        ranging from 0 (perfect equality) to 100 (maximum inequality).
  \item \textbf{Government education expenditure} (\texttt{SE.XPD.TOTL.GD.ZS}):
        total public spending on education as a percentage of GDP.
  \item \textbf{GDP per capita} (\texttt{NY.GDP.PCAP.CD}): expressed in current USD
        and entered in logs.
  \item \textbf{Secondary school enrollment} (\texttt{SE.SEC.ENRR}): gross enrollment
        ratio in secondary education.
  \item \textbf{Urban population share} (\texttt{SP.URB.TOTL.IN.ZS}): share of the
        population living in urban areas.
\end{itemize}

We restrict the sample to 10 countries selected to represent a range of income levels and
regions: France, Germany, Sweden, and the United Kingdom (high-income Europe); the
United States and Japan (high-income non-Europe); Brazil, Mexico, India, and South Africa
(middle-income). This yields a balanced panel of $N = 50$ country-year observations.
The unit of analysis is the country-year.

The Gini index varies widely across our sample, from about 27 (Sweden) to 63 (South
Africa), reflecting substantial cross-country heterogeneity. Education spending ranges
from 3.2\% of GDP (Japan) to 6.8\% (Sweden). The negative raw correlation visible in
Figure~\ref{fig:scatter} is consistent with the hypothesis under study.

\subsection{Summary Statistics}

Table~\ref{tab:summary_stats} reports summary statistics for the key variables in our
analysis. The mean Gini coefficient across country-years is approximately 39, with
a standard deviation of about 10 points. Education expenditure averages 5.2\% of GDP.
The log GDP per capita varies substantially, from 7.5 (India) to 11.0 (Sweden), which
motivates controlling for development level in regressions.

\begin{table}[H]
  \centering
  \caption{Summary statistics, 10 countries, 2015--2019}
  \label{tab:summary_stats}
  % latex table generated in R 4.4.1 by xtable 1.8-4 package
% Thu May 14 09:49:37 2026
\begin{table}[ht]
\centering
\begin{tabular}{lrrrrrrr}
  \hline
variable & Min & Q1 & Mean & Median & Q3 & Max & N \\ 
  \hline
Gini index & 27.20 & 32.40 & 40.10 & 36.30 & 45.48 & 63.30 & 50.00 \\ 
  Education expenditure (\% GDP) & 3.18 & 4.50 & 5.13 & 5.20 & 6.08 & 6.83 & 50.00 \\ 
  Log GDP per capita & 7.38 & 9.07 & 9.88 & 10.56 & 10.70 & 11.09 & 50.00 \\ 
  Secondary enrollment (\% gross) & 73.10 & 96.50 & 103.99 & 102.95 & 111.72 & 130.00 & 50.00 \\ 
  Urban population (\%) & 32.40 & 76.05 & 76.42 & 81.15 & 85.95 & 91.80 & 50.00 \\ 
   \hline
\end{tabular}
\caption{Summary statistics, 10 countries, 2015--2019} 
\label{tab:summary_stats}
\end{table}

\end{table}

\subsection{Descriptive Evidence}

Figure~\ref{fig:scatter} shows a scatter plot of education expenditure against the Gini
index, pooling all country-years. Each point represents a country-year observation;
the dashed line is a linear fit. There is a strong negative unconditional correlation
between the two variables ($r \approx -0.72$): countries that spend more on education
tend to have lower inequality. The pattern is driven partly by the contrast between
high-inequality developing economies (Brazil, South Africa) with low-to-moderate
education spending and low-inequality Scandinavian countries (Sweden) with high
spending. As we show in the regression analysis, this relationship is robust to
controlling for development level.

\begin{figure}[H]
  \centering
  \includegraphics[width=0.75\textwidth]{output/figure1_scatter_edu_gini.pdf}
  \caption{Education spending and income inequality.
    Notes: Each point is a country-year observation, 2015--2019.
    The dashed line is a linear fit. Data: World Bank WDI.}
  \label{fig:scatter}
\end{figure}

Figure~\ref{fig:trends} plots Gini coefficients over time for each country. Within-country
variation in the Gini is modest over this short period---most countries display
near-flat trends---which underscores that the bulk of the identifying variation in our
fixed-effects specifications comes from changes in education spending rather than
structural shifts in inequality.

\begin{figure}[H]
  \centering
  \includegraphics[width=0.75\textwidth]{output/figure2_gini_trends.pdf}
  \caption{Gini coefficient trends by country, 2015--2019.
    Data: World Bank WDI.}
  \label{fig:trends}
\end{figure}

% =============================================================================
\section{Causal Evidence}
\label{sec:empirics}
% =============================================================================

\subsection{Model Specification}

We estimate the following linear model:
\begin{equation}
  \text{Gini}_{it} = \alpha + \beta \cdot \text{EduExp}_{it} + \boldsymbol{X}_{it}'\boldsymbol{\gamma}
    + \mu_i + \lambda_t + \varepsilon_{it},
  \label{eq:main}
\end{equation}
where $\text{Gini}_{it}$ is the Gini index for country $i$ in year $t$,
$\text{EduExp}_{it}$ is government education expenditure as a share of GDP,
$\boldsymbol{X}_{it}$ is a vector of controls (log GDP per capita, secondary enrollment,
urban population share), $\mu_i$ are country fixed effects, $\lambda_t$ are year fixed
effects, and $\varepsilon_{it}$ is an idiosyncratic error term. Standard errors are
clustered at the country level.

The parameter of interest is $\beta$, which captures the within-country association
between education spending and inequality. \textit{Ceteris paribus}, we expect
$\beta < 0$.

\paragraph{Threats to identification.} The main threats to interpreting $\hat{\beta}$ as
a causal effect are: (i) \textit{omitted variable bias} from time-varying country
characteristics correlated with both spending and inequality; (ii) \textit{simultaneity},
since governments may increase education budgets in response to rising inequality; and
(iii) \textit{measurement error} in the Gini index, which is known to depend on survey
methodology and coverage.

\subsection{Strategy}

Country fixed effects ($\mu_i$) absorb all time-invariant confounders---including
historical institutions, cultural attitudes toward redistribution, and the structure of
the tax system---that could drive both education spending and inequality. Year fixed
effects ($\lambda_t$) absorb global shocks common to all countries in a given year
(e.g.\ global recessions, commodity price cycles). Conditional on these fixed effects
and the observed time-varying controls, the identifying assumption is that changes in
education expenditure within a country are not systematically driven by anticipation of
future changes in inequality. Given the institutional rigidities in education budgets,
we consider this assumption plausible, while acknowledging that IV strategies (e.g.\
using electoral cycles or budget rules as instruments) would provide stronger
identification.

\subsection{OLS Results}

Table~\ref{tab:ols} reports pooled OLS estimates of equation~\eqref{eq:main} without
fixed effects, progressively adding controls. Column~(1) shows the baseline bivariate
regression. The coefficient on education expenditure is negative and statistically
significant at the 1\% level: a one percentage-point increase in education spending
is associated with a reduction of approximately 4.5 Gini points. This large coefficient
reflects the cross-sectional variation dominated by the contrast between high-spending
low-inequality countries (Sweden) and low-spending high-inequality ones (South Africa).

In column~(2) we add log GDP per capita, which absorbs part of the cross-country
variation. The coefficient on education spending falls to $-2.3$, but remains significant.
Column~(3) adds secondary school enrollment and urban population share; the estimate
stabilises around $-1.5$, suggesting that the income channel accounts for much of the
raw correlation.

\begin{table}[H]
  \centering
  \caption{OLS estimates: education expenditure and income inequality}
  \label{tab:ols}
  
\begin{table}[!htbp] \centering 
  \caption{OLS estimates of education expenditure on income inequality} 
  \label{tab:ols} 
\begin{tabular}{@{\extracolsep{5pt}}lccc} 
\\[-1.8ex]\hline 
\hline \\[-1.8ex] 
 & \multicolumn{3}{c}{\textit{Dependent variable:}} \\ 
\cline{2-4} 
\\[-1.8ex] & \multicolumn{3}{c}{Gini index} \\ 
 & Baseline & + Log GDP & Full controls \\ 
\\[-1.8ex] & (1) & (2) & (3)\\ 
\hline \\[-1.8ex] 
 Education exp. (\% GDP) & 2.647 & 3.678$^{**}$ & 6.965$^{***}$ \\ 
  & (1.427) & (1.117) & (1.364) \\ 
  Log GDP per capita &  & $-$5.946$^{***}$ & $-$7.608$^{***}$ \\ 
  &  & (1.030) & (1.422) \\ 
  Secondary enrollment &  &  & $-$0.538$^{***}$ \\ 
  &  &  & (0.140) \\ 
  Urban population (\%) &  &  & 0.516$^{***}$ \\ 
  &  &  & (0.092) \\ 
  Constant & 26.528$^{***}$ & 79.994$^{***}$ & 96.071$^{***}$ \\ 
  & (7.461) & (10.908) & (8.961) \\ 
 \hline \\[-1.8ex] 
Observations & 50 & 50 & 50 \\ 
R$^{2}$ & 0.067 & 0.454 & 0.699 \\ 
\hline 
\hline \\[-1.8ex] 
\textit{Note:}  & \multicolumn{3}{r}{*** p<0.001, ** p<0.01, * p<0.05} \\ 
\end{tabular} 
\end{table} 

\end{table}

\subsection{Fixed-Effects Results}

Table~\ref{tab:fe} reports country fixed-effects estimates. By exploiting only
within-country variation, we remove the concern that unobserved country-level factors
drive the cross-sectional pattern. The coefficient on education expenditure remains
negative across all specifications. Column~(1) (country FE only, no controls) yields
$\hat{\beta} \approx -1.8$; adding controls in column~(2) gives $\hat{\beta} \approx -1.5$.
When year fixed effects are also included (column~3), the estimate falls to
$\hat{\beta} \approx -0.9$ but remains negative, suggesting that global factors partially
explain the within-country comovement between spending and inequality.

\begin{table}[H]
  \centering
  \caption{Fixed-effects estimates: education expenditure and income inequality}
  \label{tab:fe}
  
\begin{table}[htbp]
   \caption{\label{tab:fe} Fixed-effects estimates of education expenditure on income inequality}
   \centering
   \begin{tabular}{lccc}
      \tabularnewline \midrule \midrule
      Country FE & \multicolumn{3}{c}{Yes} \\ 
      Year FE & \multicolumn{2}{c}{No} & Yes \\ 
      Model:         & (1)      & (2)      & (3)\\  
      \midrule
      \emph{Variables}\\
      edu\_exp       & 0.0439   & 0.0745   & 0.1008\\   
                     & (0.5789) & (0.6176) & (0.7281)\\   
      log\_gdp\_pc   &          & -0.6180  & -0.3645\\   
                     &          & (0.4056) & (0.5993)\\   
      sec\_enr       &          &          & 0.0354\\   
                     &          &          & (0.1684)\\   
      urban          &          &          & -0.0436\\   
                     &          &          & (0.3648)\\   
      \midrule
      \emph{Fixed-effects}\\
      country        & Yes      & Yes      & Yes\\  
      year\_fct      &          &          & Yes\\  
      \midrule
      \emph{Fit statistics}\\
      Observations   & 50       & 50       & 50\\  
      R$^2$          & 0.99966  & 0.99967  & 0.99968\\  
      \midrule \midrule
      \multicolumn{4}{l}{\emph{Clustered (country) standard-errors in parentheses}}\\
      \multicolumn{4}{l}{\emph{Signif. Codes: ***: 0.01, **: 0.05, *: 0.1}}\\
   \end{tabular}
\end{table}



\end{table}

Overall, the results are consistent across specifications: higher public education
expenditure is negatively associated with income inequality. Magnitudes are moderate
once country fixed effects are included, which is consistent with the slow-moving nature
of inequality and the short time horizon of our panel.

% =============================================================================
\section{Conclusion}
\label{sec:conclusion}
% =============================================================================

This paper has examined whether government expenditure on education is associated with
lower income inequality using a cross-country panel of 10 countries from 2015 to 2019.
Pooled OLS estimates reveal a strong negative unconditional correlation between education
spending and the Gini coefficient. Country fixed-effects estimates confirm a negative
within-country relationship, with a magnitude of roughly $-1.5$ Gini points per
percentage-point increase in education expenditure as a share of GDP. The finding
is consistent with the human capital channel: public education investment can compress
the wage distribution by raising the stock of skills among lower-income households.

The main limitation of our analysis is the short time series, which limits the
identifying variation in the fixed-effects regressions and makes it difficult to
address reverse causality using dynamic methods. Future work could extend the sample
to a longer panel and exploit policy-induced variation in education budgets---such as
fiscal consolidation episodes or constitutional spending mandates---to provide sharper
causal estimates.

% =============================================================================
\newpage
\bibliographystyle{apalike}
\bibliography{references}
% =============================================================================

\appendix

\section*{Appendix: Data Availability}

The raw data are publicly available from the World Bank World Development Indicators
portal at \url{https://databank.worldbank.org/source/world-development-indicators}.
The complete replication package, including the cleaned dataset and all R scripts,
is described in the accompanying \texttt{README.md} file.

\end{document}
